% CIC (Cálculo de Construcciones Inductivas): curso exprés visual, amigable para
% TDAH e interactivo (PDF con enlaces). Sigue el Capítulo 10 de la tesis: nueve
% estaciones, un guion para explicarlo en voz alta, 18 preguntas de defensa con
% su respuesta, trampas, hoja de reglas e índice.
% Compilar dos veces:  lualatex CursoCIC.tex
\documentclass[12pt]{article}
\usepackage[a4paper,landscape,margin=11mm,bottom=17mm]{geometry}
\usepackage{amsmath}
\usepackage{fontspec}
\usepackage{unicode-math}
\usepackage[spanish,es-noshorthands]{babel}
\setmainfont{TeX Gyre Heros}
\setmonofont{DejaVu Sans Mono}[Scale=0.86]
\setmathfont{TeX Gyre DejaVu Math}
\usepackage[table]{xcolor}
\usepackage[most]{tcolorbox}
\usepackage{tikz}
\usetikzlibrary{arrows.meta,positioning,calc,shapes.geometric}
\usepackage{fontawesome5}
\usepackage{enumitem}
\usepackage{array}
\usepackage{multicol}
\usepackage{graphicx}
\usepackage{bussproofs}
\usepackage{lastpage}
\usepackage[hidelinks]{hyperref}
\pagestyle{empty}
\setlength{\parindent}{0pt}
\linespread{1.06}

% ---------- Un color por estación ----------
\definecolor{e1}{HTML}{B8901A}   % por qué no basta el CoC
\definecolor{e2}{HTML}{D9822B}   % los cuatro tiempos
\definecolor{e3}{HTML}{8C5FB8}   % Prop y Type
\definecolor{e4}{HTML}{A8456B}   % qué se puede declarar
\definecolor{e5}{HTML}{3F6FA8}   % formación e introducción
\definecolor{e6}{HTML}{2A8C8C}   % eliminación
\definecolor{e7}{HTML}{3E9A5A}   % cómputo iota
\definecolor{e8}{HTML}{2F5D8C}   % el CIC de Lean
\definecolor{e9}{HTML}{5E4A96}   % el CIC en la tesis
\definecolor{gris}{HTML}{5B6B82}
\definecolor{tldr}{HTML}{FFF4C2}
\definecolor{verde}{HTML}{3E9A5A}
\definecolor{ambar}{HTML}{E0A320}
\definecolor{rojo}{HTML}{C8463D}
\definecolor{pizarra}{HTML}{24443A}

\AtBeginDocument{\providecommand\square{\mdlgwhtsquare}}
\newcommand{\bred}{\to_\beta}
\newcommand{\bmred}{\twoheadrightarrow_\beta}
\newcommand{\beq}{=_\beta}
\newcommand{\aeq}{=_\alpha}
\newcommand{\lw}{\lambda\underline{\omega}}
\newcommand{\FV}{\mathrm{FV}}
\EnableBpAbbreviations

% ---------- Botones y navegación ----------
\newcommand{\boton}[2][gris]{\tikz[baseline=(b.base)]\node[draw=#1, fill=#1!10,
  line width=1pt, rounded corners=5pt, inner xsep=8pt, inner ysep=4pt,
  font=\bfseries\normalsize, text=#1!80!black](b){#2};}
\newcommand{\ir}[3][gris]{\hyperlink{#2}{\boton[#1]{#3}}}
\newcommand{\chip}[1]{\hyperlink{est#1}{\tikz[baseline=(c.base)]\node[fill=e#1,
  text=white, rounded corners=3pt, inner sep=2.5pt, font=\bfseries\footnotesize](c){#1};}}

% #1 destino anterior, #2 texto, #3 destino siguiente, #4 texto
\makeatletter\newcommand\pdim[1]{\strip@pt\dimexpr#1\relax}\makeatother
\newcommand{\navegacion}[4]{\AddToHookNext{shipout/foreground}{%
  \put(\pdim{0.5\paperwidth},\pdim{-\paperheight+4mm}){\makebox(0,0)[b]{%
  \ir{#1}{\faArrowLeft\ #2}\hspace{4mm}\ir{mapa}{\faMap\ Mapa}%
  \hspace{4mm}\ir{guion}{\faMicrophone\ Guion}%
  \hspace{4mm}\ir{preguntas}{\faQuestionCircle\ Preguntas}%
  \hspace{4mm}\ir{indice}{\faFont\ Índice}\hspace{4mm}\ir{#3}{#4\ \faArrowRight}}}}}

% ---------- Encabezado de estación ----------
% [#1 destino] #2 número, #3 icono, #4 título, #5 para qué sirve
\NewDocumentCommand{\estacion}{O{est#2} m m m m}{%
  \begin{tcolorbox}[colback=e#2, colframe=e#2, arc=5pt, boxrule=0pt,
      left=10pt, right=10pt, top=5pt, bottom=5pt, after skip=4pt]
    \hypertarget{#1}{}\color{white}\LARGE #3\quad\bfseries #4\hfill
    \normalfont\large\faMapSigns\ estación #2 de 9
  \end{tcolorbox}
  \begin{tcolorbox}[colback=tldr, colframe=e#2, boxrule=0.8pt, arc=4pt,
      left=10pt, right=10pt, top=4pt, bottom=4pt, after skip=6pt]
    \large\faBullseye\ \textbf{Para qué sirve:}\ #5
  \end{tcolorbox}}

% Encabezado de página general: #1 destino, #2 color, #3 icono, #4 título, #5 franja
\newcommand{\pagina}[5]{%
  \begin{tcolorbox}[colback=#2, colframe=#2, arc=5pt, boxrule=0pt,
      left=10pt, right=10pt, top=5pt, bottom=5pt, after skip=4pt]
    \hypertarget{#1}{}\color{white}\LARGE #3\quad\bfseries #4\hfill
    \normalfont\large página \thepage\ de \pageref{LastPage}
  \end{tcolorbox}
  \begin{tcolorbox}[colback=tldr, colframe=#2, boxrule=0.8pt, arc=4pt,
      left=10pt, right=10pt, top=4pt, bottom=4pt, after skip=6pt]
    \large\faBolt\ \textbf{En 10 segundos:}\ #5
  \end{tcolorbox}}

% ---------- Tarjeta ----------
% #1 color, #2 título, #3 cuerpo
\NewDocumentCommand{\tarj}{m m +m}{\begin{tcolorbox}[colback=#1!6, colframe=#1,
  boxrule=1pt, arc=4pt, left=6pt, right=6pt, top=4pt, bottom=4pt, leftrule=5pt,
  before upper=\raggedright]
  {\large\bfseries\color{#1!80!black}#2}\par\vspace{2pt}\tarjsize #3\end{tcolorbox}}
\newcommand{\tarjsize}{\normalsize}
\newcommand{\medio}{\fontsize{13}{16.2}\selectfont}
\newcommand{\grande}{\renewcommand{\tarjsize}{\large}}
\newcommand{\normal}{\renewcommand{\tarjsize}{\normalsize}}
\newcommand{\tesis}[1]{{\small\color{black!60}\faMapMarker*\ #1}}
\newcommand{\clave}[2]{\begin{tcolorbox}[colback=tldr, colframe=#1, boxrule=0.8pt,
  arc=4pt, left=8pt, right=8pt, top=3pt, bottom=3pt]\faKey\ \textbf{Clave:}\ #2\end{tcolorbox}}

% ---------- Pregunta de defensa ----------
% * = imprescindible; #2 color, #3 número, #4 pregunta, #5 respuesta, #6 dónde
\NewDocumentCommand{\preg}{s m m m +m m}{%
  \begin{tcolorbox}[colback=white, colframe=#2, boxrule=1.1pt, arc=5pt,
    left=7pt, right=7pt, top=3pt, bottom=3pt, colbacktitle=#2,
    fonttitle=\bfseries\large, before upper=\raggedright,
    title={\IfBooleanT{#1}{\faStar\ }#3.\ #4}]
    \medio\faCommentDots\ #5\par\vspace{2pt}\tesis{#6}
  \end{tcolorbox}}

% ---------- Pruébate (respuestas al revés) ----------
\newcommand{\pruebate}[3]{%
  \begin{tcolorbox}[colback=white, colframe=#1, boxrule=1.1pt, arc=5pt,
    left=8pt, right=8pt, top=4pt, bottom=4pt, sidebyside, lefthand width=0.56\linewidth,
    sidebyside gap=10pt, segmentation style={solid, #1!40}]
    {\large\bfseries\color{#1!80!black}\faBrain\ Pruébate} \ {\small(dilo en voz alta)}
    \begin{enumerate}[leftmargin=6mm, itemsep=1pt, topsep=2pt]#2\end{enumerate}
    \tcblower
    \centering{\footnotesize\color{black!50}\faSync*\ respuestas: gira la página}\par\vspace{2pt}
    \rotatebox{180}{\parbox{0.92\linewidth}{\small\raggedright #3}}
  \end{tcolorbox}}

\tikzset{
  num/.style={circle, fill=#1, text=white, font=\bfseries\Large,
    minimum size=9mm, inner sep=0pt},
  caja/.style={draw=#1, fill=#1!8, line width=1.1pt, rounded corners=5pt,
    align=center, inner sep=5pt},
}

% ---------- Código Lean ----------
% #1 color, #2 título, #3 líneas
\NewDocumentCommand{\codigo}{m m +m}{\begin{tcolorbox}[colback=black!3, colframe=#1,
  boxrule=1pt, arc=4pt, left=6pt, right=6pt, top=3pt, bottom=3pt, title={#2},
  colbacktitle=#1, fonttitle=\bfseries]\ttfamily\codesize\raggedright #3\end{tcolorbox}}
\newcommand{\codesize}{\normalsize}
\newcommand{\lean}[1]{\texttt{#1}}
\newcommand{\ired}{\to_\iota}
\newcommand{\N}{\mathbb{N}}
\newcommand{\Ind}{\mathrm{Ind}}
\newcommand{\Constr}{\mathrm{Constr}}
\newcommand{\Elim}{\mathrm{Elim}}
\newcommand{\Prop}{\mathrm{Prop}}
\newcommand{\Set}{\mathrm{Set}}
\newcommand{\Type}{\mathrm{Type}}

\begin{document}

%=====================================================================
% PORTADA
%=====================================================================
\hypertarget{portada}{}
\begin{tcolorbox}[colback=e9, colframe=e9, arc=8pt, boxrule=0pt, halign=center,
  top=14pt, bottom=14pt]
\color{white}{\fontsize{32}{38}\selectfont\bfseries \faCubes\ \ CIC: Cálculo de Construcciones Inductivas}\\[8pt]
{\LARGE Curso exprés para explicarlo en la defensa}\\[4pt]
{\large Capítulo 10 de la tesis · 9 estaciones · un guion para decirlo en voz alta · 18 preguntas con respuesta}
\end{tcolorbox}

\vspace{8pt}
\begin{minipage}[t]{0.50\textwidth}
{\Large\bfseries Cómo se usa}\\[8pt]
\begin{tikzpicture}[baseline=(current bounding box.north),
    txt/.style={align=left, text width=11.2cm, anchor=west, font=\large}]
  \node[num=gris] (a) at (0,0) {1};
  \node[txt, right=4mm of a] {Lee \textbf{el guion}: es el CIC dicho en\\ 15 segundos, en 1 minuto y en la pizarra.};
  \node[num=gris] (b) at (0,-1.55) {2};
  \node[txt, right=4mm of b] {Recorre \textbf{el mapa}: 9 estaciones, de\\ «¿por qué no basta el CoC?» a tu tesis.};
  \node[num=gris] (c) at (0,-3.1) {3};
  \node[txt, right=4mm of c] {En cada estación: \textbf{para qué sirve} $\to$\\ \textbf{tarjetas} $\to$ \textbf{pruébate} (respuestas al revés).};
  \node[num=gris] (d) at (0,-4.65) {4};
  \node[txt, right=4mm of d] {Termina con las \textbf{18 preguntas}:\\ tapa la respuesta y contesta en voz alta.};
\end{tikzpicture}
\end{minipage}\hfill
\begin{minipage}[t]{0.46\textwidth}
{\Large\bfseries Empieza por aquí (haz clic)}\\[10pt]
\begin{tcolorbox}[colback=e2, colframe=e2, arc=6pt, halign=center, hyperlink=guion]
  \color{white}\Large\faMicrophone\ \textbf{El guion}\\ \normalsize cómo explicar el CIC en voz alta\end{tcolorbox}
\begin{tcolorbox}[colback=rojo, colframe=rojo, arc=6pt, halign=center, hyperlink=preguntas]
  \color{white}\Large\faQuestionCircle\ \textbf{Las 18 preguntas}\\ \normalsize lo que te van a preguntar, con la respuesta\end{tcolorbox}
\begin{tcbraster}[raster columns=3, raster column skip=6pt, raster equal height]
\begin{tcolorbox}[colback=e5, colframe=e5, arc=6pt, halign=center, hyperlink=mapa]
  \color{white}\large\faMap\ \textbf{Mapa}\\ \small 9 estaciones\end{tcolorbox}
\begin{tcolorbox}[colback=gris, colframe=gris, arc=6pt, halign=center, hyperlink=reglas]
  \color{white}\large\faScroll\ \textbf{Reglas}\\ \small en una página\end{tcolorbox}
\begin{tcolorbox}[colback=gris, colframe=gris, arc=6pt, halign=center, hyperlink=indice]
  \color{white}\large\faFont\ \textbf{Índice}\\ \small A–Z\end{tcolorbox}
\end{tcbraster}
\end{minipage}

\vfill
\begin{tcolorbox}[colback=tldr, colframe=gris, arc=4pt, boxrule=0.8pt, halign=center]
\large\faLightbulb\ \textbf{La frase que lo une todo:}\ \ CIC $=$ CoC $+$ tipos inductivos.\ \ Un tipo
inductivo se declara por sus \textbf{constructores},\\ y el sistema entrega su \textbf{recursor}:
con él se define por recursión y se demuestra por inducción.
\end{tcolorbox}

\newpage
%=====================================================================
% MAPA
%=====================================================================
\pagina{mapa}{gris}{\faMap}{El mapa: 9 estaciones}{el CoC ya tiene funciones y «para todo».
El CIC le añade tipos de datos con inducción. Eso es lo que usa Lean, y lo que usa tu tesis.}

\newcommand{\parada}[4]{\begin{tcolorbox}[colback=e#1!8, colframe=e#1, boxrule=1.2pt,
  arc=6pt, halign=center, left=3pt, right=3pt, top=3pt, bottom=3pt, hyperlink=est#1]
  \tikz\node[num=e#1]{#1};\\[2pt]{\large\bfseries\color{e#1!80!black}#2}\\[1pt]
  {\small #3}\\[1pt]{\footnotesize\color{black!55}#4}\end{tcolorbox}}

\begin{tcbraster}[raster columns=5, raster column skip=6pt, raster row skip=6pt,
  raster equal height=rows]
\parada{1}{¿Por qué no basta el CoC?}{datos codificados con funciones}{§10.2}
\parada{2}{Los cuatro tiempos}{formar, introducir, eliminar, computar}{§10.3 · Nota 10.6}
\parada{3}{\lean{Prop} y \lean{Type}}{proposiciones y datos}{Nota 10.2}
\parada{4}{Qué se puede declarar}{aridad, constructor, positividad}{§10.3.2}
\parada{5}{Formación e introducción}{las reglas Ind y Constr}{§10.3.3}
\parada{6}{Eliminación}{el recursor: recursión e inducción}{§10.3.1 y §10.3.4}
\parada{7}{Cómputo $\iota$}{el recursor frente a un constructor}{§10.3.5}
\parada{8}{El CIC de Lean}{tres diferencias y las tácticas}{Nota 10.7 · §10.4}
\parada{9}{El CIC en mi tesis}{\lean{Formula}, \lean{ND}, listas}{Capítulos 11 y 12}
\begin{tcolorbox}[colback=e2!10, colframe=e2, boxrule=1.2pt, arc=6pt, halign=center,
  left=3pt, right=3pt, top=3pt, bottom=3pt, hyperlink=guion]
  \tikz\node[num=e2]{\faMicrophone};\\[2pt]{\large\bfseries\color{e2!80!black}El guion}\\[1pt]
  {\small todo lo anterior, dicho en voz alta}\\[1pt]{\footnotesize\color{black!55}empieza aquí}\end{tcolorbox}
\end{tcbraster}

\vspace{4pt}
\begin{center}
\begin{tikzpicture}[font=\normalsize, >={Stealth[length=3mm]},
  paso/.style={draw=#1, fill=#1!10, line width=1.1pt, rounded corners=5pt,
    align=center, text width=3.6cm, minimum height=1.45cm, inner sep=4pt}]
  \node[paso=e1] (a) {\textbf{CoC}\\ \small solo $\lambda$ y $\Pi$:\\ los datos se codifican};
  \node[paso=e1, right=7mm of a] (b) {\textbf{El problema}\\ \small hay iteración,\\ falta la inducción};
  \node[paso=e2, right=7mm of b] (c) {\textbf{Tipo inductivo}\\ \small se declaran\\ los constructores};
  \node[paso=e6, right=7mm of c] (d) {\textbf{Recursor $+$ $\iota$}\\ \small recursión, inducción\\ y cómputo};
  \node[paso=e8, right=7mm of d] (e) {\textbf{Lean}\\ \small una variante del CIC\\ en el núcleo};
  \node[paso=e9, right=7mm of e] (f) {\textbf{Mi tesis}\\ \small \lean{Formula} y \lean{ND}\\ son inductivos};
  \foreach \u/\v in {a/b,b/c,c/d,d/e,e/f} \draw[->, line width=1.3pt, black!60] (\u) -- (\v);
\end{tikzpicture}
\end{center}

\begin{tcbraster}[raster columns=2, raster column skip=8pt, raster equal height]
\begin{tcolorbox}[colback=tldr, colframe=black!50, arc=5pt, boxrule=0.8pt]
\faRunning\ \textbf{Ruta rápida (15 minutos):}\ \ir[e2]{guion}{\faMicrophone\ Guion} $\to$ \chip{2} $\to$ \chip{6} $\to$
\chip{7} $\to$ \chip{9} $\to$ \ir[rojo]{preguntas}{\faQuestionCircle\ Preguntas}
\end{tcolorbox}
\begin{tcolorbox}[colback=tldr, colframe=black!50, arc=5pt, boxrule=0.8pt]
\faExclamationTriangle\ \textbf{Recuerda:} el Capítulo 10 es \textbf{exposición}. Las reglas son de
Paulin-Mohring (1993). Tú no demuestras nada sobre el CIC: lo \textbf{usas}.
\end{tcolorbox}
\end{tcbraster}
\navegacion{portada}{Portada}{guion}{Guion}

\newpage
%=====================================================================
% GUION
%=====================================================================
\pagina{guion}{e2}{\faMicrophone}{El guion: explica el CIC en voz alta}{elige la versión según el
tiempo que te den. Las tres dicen lo mismo; solo cambia el detalle.}

\begin{minipage}[t]{0.37\textwidth}\vspace{0pt}
\begin{tcolorbox}[colback=e2!8, colframe=e2, boxrule=1.3pt, arc=6pt, left=7pt, right=7pt,
  title={\faStopwatch\ En 15 segundos}, colbacktitle=e2, fonttitle=\bfseries\large,
  before upper=\raggedright]
\large «El CIC es el Cálculo de Construcciones más \textbf{tipos inductivos}. Un tipo inductivo se
declara dando sus \textbf{constructores}, y el sistema genera su \textbf{eliminador}, que sirve a la
vez para la recursión y para la inducción. Con variantes, es la teoría que implementa el núcleo de Lean.»
\end{tcolorbox}
\begin{tcolorbox}[colback=pizarra, colframe=black!70, boxrule=2pt, arc=4pt, left=8pt, right=8pt,
  top=5pt, bottom=5pt, title={\faChalkboard\ En la pizarra: escribe solo esto},
  colbacktitle=black!70, fonttitle=\bfseries\large]
\color{white}\ttfamily\normalsize\raggedright
\newcommand{\eti}[2]{\hfill{\rmfamily\small\color{ambar}\tikz[baseline=(q.base)]\node[circle,fill=ambar,text=black,inner sep=1.2pt,font=\bfseries\footnotesize](q){#1};\ #2}}
inductive Nat\eti{1}{formar}\\
| zero : Nat\eti{2}{introducir}\\
| succ : Nat → Nat\\[5pt]
Nat.rec :\ \ P 0 →\eti{3}{eliminar}\\
\ \ (∀ n, P n → P (n+1)) →\\
\ \ ∀ n, P n\\[5pt]
rec z s 0\ \ \ \ \ $\longrightarrow$\ z\eti{4}{computar ($\iota$)}\\
rec z s (n+1) $\longrightarrow$\ s n (rec z s n)
\end{tcolorbox}
\end{minipage}\hfill
\begin{minipage}[t]{0.61\textwidth}\vspace{0pt}
\begin{tcolorbox}[colback=white, colframe=e2, boxrule=1.3pt, arc=6pt, left=7pt, right=7pt,
  title={\faClock\ En 1 minuto: seis frases, en este orden}, colbacktitle=e2,
  fonttitle=\bfseries\large, before upper=\raggedright]
\large\renewcommand{\arraystretch}{1.25}
\begin{tabular}{@{}c>{\raggedright\arraybackslash}p{14.1cm}@{}}
\tikz[baseline=-4pt]\node[num=e1, minimum size=7mm, font=\bfseries\large]{1}; &
«El CoC ya tiene funciones y cuantificación. Por Curry–Howard, las proposiciones son tipos y las
pruebas son términos.»\\
\tikz[baseline=-4pt]\node[num=e1, minimum size=7mm, font=\bfseries\large]{2}; &
«Pero en el CoC puro los datos se \textbf{codifican} con funciones. Eso da iteración; el principio de
inducción no sale de la codificación.»\\
\tikz[baseline=-4pt]\node[num=e2, minimum size=7mm, font=\bfseries\large]{3}; &
«El CIC añade \textbf{tipos inductivos}: se declaran los constructores. Para los naturales, el cero y el sucesor.»\\
\tikz[baseline=-4pt]\node[num=e6, minimum size=7mm, font=\bfseries\large]{4}; &
«De esa declaración salen cuatro cosas: el tipo, sus constructores, su \textbf{recursor} y la regla de
cómputo $\iota$.»\\
\tikz[baseline=-4pt]\node[num=e6, minimum size=7mm, font=\bfseries\large]{5}; &
«El recursor es una sola pieza con dos usos: si devuelve datos, es la \textbf{recursión}; si devuelve
pruebas, su tipo es el \textbf{principio de inducción}.»\\
\tikz[baseline=-4pt]\node[num=e9, minimum size=7mm, font=\bfseries\large]{6}; &
«En mi tesis, las fórmulas y la derivabilidad \lean{ND} son tipos inductivos. “Inducción sobre la
derivación” es aplicar el recursor de \lean{ND}.»
\end{tabular}
\end{tcolorbox}
\end{minipage}

\vfill
\begin{tcolorbox}[colback=tldr, colframe=black!50, boxrule=0.8pt, arc=5pt, before upper=\raggedright]
\faShield*\ \textbf{Si preguntan «¿eso lo demostró usted?»:} «No: el Capítulo 10 es exposición y sigo a
Paulin-Mohring (1993). Mi aporte está en los Capítulos 11 y 12.»
\end{tcolorbox}
\navegacion{mapa}{Mapa}{est1}{Estación 1}

\newpage
%=====================================================================
% ESTACIÓN 1: POR QUÉ NO BASTA EL COC
%=====================================================================
\grande
\estacion{1}{\faPuzzlePiece}{¿Por qué no basta el CoC?}{es la motivación del capítulo. Si
explicas bien este problema, los tipos inductivos se entienden como su solución.}

\begin{minipage}[t]{0.49\textwidth}\vspace{0pt}
\tarj{e1}{\faCode\ Codificación impredicativa}{
En el CoC puro no hay naturales, listas ni conectivos: todo se construye con $\lambda$ y $\Pi$.\par\vspace{3pt}
\centerline{$\mathrm{Nat} := \forall A,\ A \to (A\to A)\to A$}\par\vspace{3pt}
$0=\lambda A\,x\,f.\,x \qquad 2=\lambda A\,x\,f.\,f\,(f\,x)$\par\vspace{2pt}
El número $n$ \textbf{es} «aplicar $f$ un total de $n$ veces».}
\normal\tarj{gris}{\faTable\ Más codificaciones (Tabla 10.1)}{
\renewcommand{\arraystretch}{1.2}
\begin{tabular}{@{}ll@{}}
Falsedad $\bot$ & $\forall A:\Prop,\ A$\\
Par $A\times B$ & $\forall C:\Prop,\ (A\to B\to C)\to C$\\
Igualdad $x=y$ & $\forall P:A\to\Prop,\ P\,y\to P\,x$\quad (Leibniz)
\end{tabular}\\[2pt]
\textbf{Impredicativa:} cuantifica sobre una colección que incluye lo que se está definiendo.}\grande
\end{minipage}\hfill
\begin{minipage}[t]{0.49\textwidth}\vspace{0pt}
\tarj{rojo}{\faExclamationTriangle\ Tres límites de la codificación}{
\textbf{1. Solo itera.} No hay acceso directo a $n-1$: el predecesor se calcula con pares
y cuesta $n$ pasos.\par\vspace{3pt}
\textbf{2. La inducción no sale.} De la codificación no se obtiene\par
\centerline{$\forall P,\ P\,0\to(\forall n,\ P\,n\to P\,(S\,n))\to\forall n,\ P\,n$}\par
\textbf{Citado:} Geuvers (2001).\par\vspace{3pt}
\textbf{3. En Lean no son datos.} Lo cuantificado sobre \lean{Prop} vive en \lean{Prop}, y ahí
dos pruebas no se distinguen.}
\clave{e1}{Tener $0$ y $S$ no basta. Falta un principio para razonar sobre
\textbf{cualquier} habitante del tipo. Eso es lo que añade el CIC.}
\tesis{Codificaciones: §10.2.2 · Iteración, costo e inducción: §10.2.3}
\end{minipage}

\vfill
\pruebate{e1}{
\item ¿Cómo se escribe el número 2 en la codificación de Church?
\item ¿Qué da la codificación y qué no da?
\item ¿Quién prueba que la inducción no es derivable? ¿Lo demuestras tú?
}{\textbf{1.} $\lambda A\,x\,f.\,f\,(f\,x)$: aplicar $f$ dos veces.\quad
\textbf{2.} Da iteración. No da la inducción dependiente ni un predecesor barato.\quad
\textbf{3.} Geuvers (2001). No: lo cito.}
\navegacion{guion}{Guion}{est2}{Estación 2}

\newpage
%=====================================================================
% ESTACIÓN 2: LOS CUATRO TIEMPOS
%=====================================================================
\grande
\estacion{2}{\faLayerGroup}{Los cuatro tiempos de un tipo inductivo}{es el molde de todo el
capítulo. Cualquier tipo inductivo, también \lean{Formula} y \lean{ND}, se explica con estas cuatro piezas.}

\begin{center}
\begin{tikzpicture}[>={Stealth[length=3mm]},
  pieza/.style={draw=#1, fill=#1!9, line width=1.3pt, rounded corners=6pt, align=center,
    text width=6.0cm, minimum height=3.0cm, inner sep=5pt, font=\normalsize}]
  \node[draw=e2, fill=white, line width=1.6pt, rounded corners=6pt, inner sep=7pt,
    font=\ttfamily\large, align=left] (decl) at (0,0)
    {inductive Nat\ \ |\ \ zero : Nat\ \ |\ \ succ : Nat → Nat};
  \node[pieza=e5, anchor=north] (f) at (-10.2,-1.75)
    {{\large\bfseries\color{e5!80!black}1 · Formación}\\[2pt] \emph{¿Qué tipo es?}\\[3pt]
     \lean{Nat : Type}\\[3pt] {\small regla \textbf{Ind}}};
  \node[pieza=e5, anchor=north] (i) at (-3.4,-1.75)
    {{\large\bfseries\color{e5!80!black}2 · Introducción}\\[2pt] \emph{¿Cómo construyo uno?}\\[3pt]
     \lean{Nat.zero},\ \ \lean{Nat.succ}\\[3pt] {\small las \textbf{únicas} maneras · \textbf{Constr}}};
  \node[pieza=e6, anchor=north] (e) at (3.4,-1.75)
    {{\large\bfseries\color{e6!80!black}3 · Eliminación}\\[2pt] \emph{¿Cómo uso uno?}\\[3pt]
     \lean{Nat.rec}\\[3pt] {\small la \textbf{única} manera · \textbf{Elim}}};
  \node[pieza=e7, anchor=north] (c) at (10.2,-1.75)
    {{\large\bfseries\color{e7!80!black}4 · Cómputo}\\[2pt] \emph{¿Y si elimino lo que construí?}\\[3pt]
     reducción $\iota$\\[3pt] {\small recursor frente a constructor}};
  \foreach \n in {f,i,e,c} \draw[->, line width=1.3pt, black!55] (decl.south) -- (\n.north);
\end{tikzpicture}
\end{center}

\normal\begin{tcbraster}[raster columns=3, raster column skip=7pt, raster equal height]
\tarj{e2}{\faSeedling\ La menor colección cerrada}{
$\N$ es la \textbf{menor} colección que contiene $0$ y está cerrada bajo $S$. No hay más
habitantes que los que producen los constructores. Por eso vale la inducción: si algo vale
para $0$ y se conserva con $S$, vale para todos.}
\tarj{e2}{\faProjectDiagram\ El mismo molde en la lógica}{
En Lean los conectivos \textbf{no son primitivos}: son tipos inductivos (Listado 10.7).\\[2pt]
\lean{And}: un constructor, \lean{intro}.\\
\lean{Or}: dos, \lean{inl} e \lean{inr}.\\
\lean{False}: \textbf{ninguno}; por eso de \lean{False} se sigue todo.}
\tarj{e2}{\faHistory\ De dónde viene}{
Coquand y Paulin (1990) extienden el CoC con definiciones inductivas primitivas.
Las reglas están en Paulin-Mohring (1993) y la metateoría en Werner (1994).
La teoría de Martin-Löf (1984) usa la misma idea: constructores con reglas de eliminación.}
\end{tcbraster}

\vfill
\pruebate{e2}{
\item Di los cuatro tiempos en orden.
\item ¿Qué significa que los constructores sean las «únicas maneras»?
\item ¿Cuántos constructores tiene \lean{False}? ¿Qué se sigue de eso?
}{\textbf{1.} Formación, introducción, eliminación, cómputo.\quad
\textbf{2.} Que todo habitante se obtiene con ellos: el tipo es la menor colección cerrada.\quad
\textbf{3.} Ninguno. Su eliminador no tiene casos que probar: de \lean{False} se sigue cualquier cosa.}
\navegacion{est1}{Estación 1}{est3}{Estación 3}

\newpage
%=====================================================================
% ESTACIÓN 3: PROP Y TYPE
%=====================================================================
\grande
\estacion{3}{\faBalanceScale}{\lean{Prop} y \lean{Type}: proposiciones y datos}{explica una decisión
de tu código: \lean{Formula} vive en \lean{Type} y \lean{ND} vive en \lean{Prop}.}

\begin{tcolorbox}[colback=white, colframe=e3, boxrule=1.5pt, arc=6pt, halign=center,
  top=4pt, bottom=4pt]
{\Large $\underbrace{0}_{\text{\normalsize dato}} :\ \underbrace{\N}_{\text{\normalsize tipo}} :\
\underbrace{\texttt{Type}}_{\text{\normalsize sort}} :\ \texttt{Type 1} :\ \cdots
\qquad\qquad
\underbrace{h}_{\text{\normalsize prueba}} :\ \underbrace{p}_{\text{\normalsize proposición}} :\
\underbrace{\texttt{Prop}}_{\text{\normalsize sort}} :\ \texttt{Type}$}\\[3pt]
{\normalsize Un \textbf{sort} es un «tipo de tipos»: así como un número pertenece a $\N$, un tipo pertenece a un sort.}
\end{tcolorbox}

\begin{minipage}[t]{0.49\textwidth}\vspace{0pt}
\tarj{e3}{\faGavel\ \lean{Prop}: proposiciones}{
Sus habitantes son \textbf{pruebas}. Importa que haya una, no cuál: son irrelevantes para el cómputo.\par\vspace{2pt}
\textbf{Impredicativa:} $\forall p:\texttt{Prop},\ p$ cuantifica sobre todas las proposiciones y
sigue siendo una proposición.}
\tarj{e3}{\faDatabase\ \lean{Type}: datos}{
Sus habitantes son \textbf{datos} que sí se distinguen: $0\neq 1$.\par\vspace{2pt}
\textbf{Predicativa:} cuantificar sobre un universo sube de universo.
$\N\to\texttt{Type}$ vive en \lean{Type 1}.}
\tesis{Nota 10.2 · GTS: §10.3.1, tras la Nota 10.1}
\end{minipage}\hfill
\begin{minipage}[t]{0.49\textwidth}\vspace{0pt}
\tarj{e9}{\faThumbtack\ En mi tesis}{
\lean{Formula : Type}\quad es sintaxis: un \textbf{dato}.\\
\lean{ND Γ φ : Prop}\quad es una afirmación: «$\varphi$ se deriva de $\Gamma$».}
\normal\tarj{gris}{\faSitemap\ Sistemas de tipos generalizados (GTS o PTS)}{
Tres datos $(\mathcal S,\mathcal A,\mathcal R)$: \textbf{sorts}, \textbf{axiomas} (qué sort habita
cuál) y \textbf{reglas} (qué productos se permiten).\par\vspace{2pt}
CoC puro:\ \ $\mathcal S=\{\Set,\Type_S\}$,\ \ $\mathcal A=\{\Set:\Type_S\}$,\ \ $\mathcal R=\mathcal S\times\mathcal S$.}
\tarj{gris}{\faExchange*\ Coq y Lean}{
Coq: \lean{Prop}, \lean{Set} y \lean{Type}.\quad Lean \textbf{elimina} \lean{Set}: queda una \lean{Prop}
impredicativa y la jerarquía predicativa \lean{Type u}.}\grande
\end{minipage}

\vfill
\pruebate{e3}{
\item ¿Qué es un sort?
\item ¿Qué significa que \lean{Prop} sea impredicativa?
\item ¿Por qué \lean{Formula} está en \lean{Type} y \lean{ND} en \lean{Prop}?
}{\textbf{1.} Un tipo de tipos: \lean{Prop}, \lean{Type}, \lean{Type 1}\ldots\quad
\textbf{2.} Que cuantificar sobre todas las proposiciones da otra proposición.\quad
\textbf{3.} Una fórmula es un dato que se inspecciona; una derivación es una prueba: solo importa que exista.}
\navegacion{est2}{Estación 2}{est4}{Estación 4}

\newpage
%=====================================================================
% ESTACIÓN 4: QUÉ SE PUEDE DECLARAR
%=====================================================================
\grande
\estacion{4}{\faClipboardCheck}{Qué se puede declarar: aridad, constructor, positividad}{no toda declaración se acepta: tres condiciones mantienen consistente el sistema.}

\begin{minipage}[t]{0.49\textwidth}\vspace{0pt}
\tarj{e4}{\faRulerHorizontal\ Aridad: argumentos y un sort final}{
\centerline{$Ar ::= s \ \mid\ (x:M)\,Ar$}\par\vspace{2pt}
\lean{nat}: $\Set$.\quad \lean{list}: $(A{:}\Set)\,\Set$.\quad En Lean se llama \textbf{firma}.}
\tarj{e4}{\faSlidersH\ Parámetro o índice}{
\textbf{Parámetro:} fijo en toda la declaración, como el $\alpha$ de \lean{List α}.\par\vspace{2pt}
\textbf{Índice:} lo fija cada constructor en su conclusión. En \lean{ND}, $\Gamma$ y $\varphi$ son
índices: es una \textbf{familia inductiva} (Dybjer, 1994).}
\normal\codigo{rojo}{\faLaptopCode\ Extra, para mostrar en vivo: Lean lo rechaza}{
inductive Malo : Type\\
| mk : (Malo → False) → Malo\\[2pt]
{\small\color{rojo!80!black}(kernel) arg \#1 of 'Malo.mk' has a non positive\\ occurrence of the datatypes being declared}}\grande
\end{minipage}\hfill
\begin{minipage}[t]{0.49\textwidth}\vspace{0pt}
\tarj{e4}{\faHammer\ Forma de constructor}{
Un tipo cuya \textbf{conclusión} es $X$, el tipo que se define, y donde $X$ solo aparece de forma
estrictamente positiva en los argumentos.\\ Para \lean{Nat}:\ \ $X$\ \ (cero)\ \ y\ \ $X\to X$\ \ (sucesor).}
\tarj{rojo}{\faBan\ Positividad estricta}{
$X$ \textbf{no puede aparecer a la izquierda de una flecha}, a ninguna profundidad.\par\vspace{2pt}
En $(X\to A)\to B$ la $X$ es covariante, pero \textbf{no} es estrictamente positiva.\par\vspace{2pt}
\textbf{Por qué:} con un constructor $(X\to\bot)\to X$ se construye un punto fijo y un habitante
de $\bot$.}
\tesis{§10.3.2 · el ejemplo \lean{Malo} no está en la tesis: lo comprobé en Lean}
\end{minipage}

\vfill
\pruebate{e4}{
\item ¿Qué diferencia hay entre un parámetro y un índice?
\item ¿Es $X$ estrictamente positiva en $(X\to A)\to B$?
\item ¿Qué pasaría si se admitiera el constructor $(X\to\bot)\to X$?
}{\textbf{1.} El parámetro es el mismo en toda la declaración; el índice lo decide cada constructor.\quad
\textbf{2.} No: está a la izquierda de una flecha, aunque sea covariante.\quad
\textbf{3.} Se obtiene un habitante de $\bot$: el sistema sería inconsistente.}
\navegacion{est3}{Estación 3}{est5}{Estación 5}

\newpage
%=====================================================================
% ESTACIÓN 5: FORMACIÓN E INTRODUCCIÓN
%=====================================================================
\grande
\estacion{5}{\faIndustry}{Formación e introducción: las reglas Ind y Constr}{son las dos
reglas que dicen qué es declarar un tipo inductivo y qué tipo tienen sus constructores.}

\begin{tcbraster}[raster columns=2, raster column skip=8pt, raster equal height]
\begin{tcolorbox}[colback=white, colframe=e5, boxrule=1.5pt, arc=6pt, halign=center,
  title={Formación (Ind)}, colbacktitle=e5, fonttitle=\bfseries\large]
{\large $\displaystyle\frac{\Gamma,\,X:A\ \vdash\ C_i:s\qquad (i=1,\dots,n)}
{\Gamma\ \vdash\ \Ind(X:A)\{C_1\mid\dots\mid C_n\}\ :\ A}$}\\[5pt]
{\normalsize si $A$ es una aridad de $s$ y cada $C_i$ es una forma de constructor de $X$}
\end{tcolorbox}
\begin{tcolorbox}[colback=white, colframe=e5, boxrule=1.5pt, arc=6pt, halign=center,
  title={Introducción (Constr)}, colbacktitle=e5, fonttitle=\bfseries\large]
{\large $\displaystyle\frac{\Gamma\ \vdash\ I:A}
{\Gamma\ \vdash\ \Constr(i,I)\ :\ C_i[X:=I]}$}\\[5pt]
{\normalsize donde $I=\Ind(X:A)\{C_1\mid\dots\mid C_n\}$ y $1\le i\le n$}
\end{tcolorbox}
\end{tcbraster}

\begin{minipage}[t]{0.49\textwidth}\vspace{0pt}
\tarj{e5}{\faBookReader\ Cómo se lee}{
$X$ nombra al tipo \textbf{dentro} de su propia definición; la barra $\mid$ separa los constructores.
$\Constr(i,I)$ es el constructor $i$: su tipo es $C_i$ con $X:=I$.}
\tarj{e5}{\faCalculator\ El ejemplo: los naturales (Ejemplo 10.3)}{
\centerline{$\texttt{nat}\ \equiv\ \Ind(X:\Set)\{\,X \mid X\to X\,\}$}\par\vspace{3pt}
$\Constr(1,\texttt{nat}):\texttt{nat}$\hfill el cero\\
$\Constr(2,\texttt{nat}):\texttt{nat}\to\texttt{nat}$\hfill el sucesor}
\tesis{§10.3.3 · Ejemplo 10.3}
\end{minipage}\hfill
\begin{minipage}[t]{0.49\textwidth}\vspace{0pt}
\tarj{rojo}{\faExclamationTriangle\ No es sintaxis de Lean}{
$\Ind$, $\Constr$ y $\Elim$ son \textbf{pseudo-términos} de Paulin-Mohring. En Lean se escribe
\lean{inductive}, los nombres de los constructores y \lean{.rec}.}
\tarj{gris}{\faGripLinesVertical\ Tres usos del torniquete $\vdash$}{
$\Gamma\vdash t:A$\hfill \textbf{tipado}: $t$ tiene tipo $A$\\
$\Gamma\vdash\varphi$\hfill \textbf{derivabilidad}: mi \lean{ND}\\
$\vdash\varphi$ en Lean\hfill la \textbf{meta} pendiente}
\end{minipage}

\vfill
\pruebate{e5}{
\item Escribe los naturales con $\Ind$.
\item ¿Qué tipo tiene $\Constr(2,\texttt{nat})$? ¿De dónde sale?
\item ¿Se escribe $\Ind$ en un archivo de Lean?
}{\textbf{1.} $\Ind(X:\Set)\{X\mid X\to X\}$.\quad
\textbf{2.} $\texttt{nat}\to\texttt{nat}$: es $X\to X$ con $X:=\texttt{nat}$.\quad
\textbf{3.} No. Es notación de Paulin-Mohring; en Lean se escribe \lean{inductive}.}
\navegacion{est4}{Estación 4}{est6}{Estación 6}

\newpage
%=====================================================================
% ESTACIÓN 6: ELIMINACIÓN
%=====================================================================
\grande
\estacion{6}{\faRecycle}{Eliminación: el recursor}{es la pieza central. De aquí salen la recursión
y la inducción. «Por inducción sobre la derivación» es aplicar un recursor.}

\begin{tcolorbox}[colback=white, colframe=e6, boxrule=1.6pt, arc=6pt, halign=center,
  top=3pt, bottom=3pt]
{\Large $\texttt{Nat.rec}\ :\
\underbrace{(P:\N\to\texttt{Sort}\ u)}_{\text{\normalsize\textbf{motivo}: lo que quiero}}\ \to\
\underbrace{P\,0}_{\text{\normalsize caso base}}\ \to\
\underbrace{(\forall n,\ P\,n\to P\,(n{+}1))}_{\text{\normalsize paso, con hipótesis inductiva}}\ \to\
\underbrace{\forall n,\ P\,n}_{\text{\normalsize conclusión}}$}\\[2pt]
{\normalsize Con \mbox{$u=0$}, que es \lean{Prop}, el recursor es \textbf{inducción}. Con \mbox{$u\ge 1$} es \textbf{recursión}.}
\end{tcolorbox}

\begin{tcbraster}[raster columns=2, raster column skip=8pt, raster equal height]
\tarj{e6}{\faRedo\ No dependiente: recursión}{
El motivo $Q$ es un tipo \textbf{fijo}. Sirve para \textbf{definir funciones}.\\[2pt]
Manejadores:\ \ cero: $Q$\qquad sucesor: $\N\to Q\to Q$\\[2pt]
Es la recursión primitiva:\ \ $g(0)=x$,\ \ $g(n{+}1)=f(n,g(n))$.}
\tarj{e6}{\faStream\ Dependiente: inducción}{
El motivo $Q$ \textbf{depende del valor}. Sirve para \textbf{demostrar}.\\[2pt]
Manejadores:\ \ $Q\,0$\quad y\quad $(n{:}\N)\to Q\,n\to Q\,(S\,n)$\\[2pt]
El segundo es el \textbf{paso inductivo}. Peano lo postula; el CIC lo genera.}
\end{tcbraster}

\tarj{rojo}{\faLock\ Restricción: de \lean{Prop} solo hacia \lean{Prop} (Nota 10.5)}{
Un inductivo que vive en \lean{Prop} solo se elimina hacia \lean{Prop}, salvo los casos \textbf{vacío} y
\textbf{unitario} (eliminación singleton). \textbf{Por qué:} la eliminación fuerte sobre una sort
impredicativa haría inconsistente el sistema, y las pruebas dejarían de ser borrables.
\textbf{En Lean:} el motivo de \lean{Or.rec} y el de \lean{ND.rec} llegan a \lean{Prop}; el de
\lean{False.rec} y el de \lean{Nat.rec}, a \lean{Sort u}.}
\tesis{Recursor de $\N$: §10.3.1 y Nota 10.1 · Reglas generales: §10.3.4 y \hyperlink{reglas}{hoja de reglas} · Ejemplo 10.4}

\vfill
\pruebate{e6}{
\item ¿Qué es el motivo?
\item ¿Qué cambia entre la eliminación no dependiente y la dependiente?
\item ¿Puedo definir una función a \lean{Nat} por casos sobre una prueba de $p\lor q$?
}{\textbf{1.} El tipo de retorno del recursor: la propiedad que se prueba o el tipo que se calcula.\quad
\textbf{2.} En la dependiente el motivo depende del valor: de recursión se pasa a inducción.\quad
\textbf{3.} No: \lean{Or} vive en \lean{Prop} y solo se elimina hacia \lean{Prop}.}
\navegacion{est5}{Estación 5}{est7}{Estación 7}

\newpage
%=====================================================================
% ESTACIÓN 7: CÓMPUTO IOTA
%=====================================================================
\grande
\estacion{7}{\faCogs}{Cómputo: la reducción $\iota$}{es lo que hace que una definición recursiva \textbf{calcule}. Detrás de cada \lean{rfl} está esto.}

\begin{tcolorbox}[colback=white, colframe=e7, boxrule=1.6pt, arc=6pt, halign=center,
  top=5pt, bottom=5pt]
{\Large $\Elim(\texttt{zero},Q)\{f_0\mid f_S\}\ \ired\ f_0
\qquad\qquad
\Elim(S\,n,Q)\{f_0\mid f_S\}\ \ired\ f_S\ n\ \bigl(\Elim(n,Q)\{f_0\mid f_S\}\bigr)$}\\[4pt]
{\normalsize\textbf{En una frase:} el recursor frente a un constructor elige el manejador de ese
constructor, le pasa sus argumentos y repite en los argumentos recursivos.}
\end{tcolorbox}

\begin{minipage}[t]{0.49\textwidth}\vspace{0pt}
\normal\codigo{e7}{\faLaptopCode\ En Lean: las dos reglas valen por \texttt{rfl} (Listado 10.4)}{
example : Nat.rec z s 0 = z := rfl\\[3pt]
example : Nat.rec z s (n+1)\\
\ \ \ \ \ \ \ \ \ \ = s n (Nat.rec z s n) := rfl}
\tarj{e7}{\faShoePrints\ Un cálculo paso a paso}{
La suma recurre sobre el \textbf{segundo} argumento. Al desplegarla ($\delta$) y reducir ($\beta$, $\iota$):\\
$m+0\ \longrightarrow\ m$\qquad $m+S\,n\ \longrightarrow\ S\,(m+n)$\par\vspace{3pt}
\centerline{$2+1\ =\ 2+S\,0\ \longrightarrow\ S\,(2+0)\ \longrightarrow\ S\,2\ =\ 3$}}
\grande\tesis{§10.3.5 · §10.4.2 · regla general: \hyperlink{reglas}{hoja de reglas}}
\end{minipage}\hfill
\begin{minipage}[t]{0.49\textwidth}\vspace{0pt}
\tarj{e7}{\faEquals\ Dos igualdades distintas}{
\textbf{Por definición.} El núcleo reduce los dos lados y compara; se cierra con \lean{rfl}.\\
\quad$\beta$ aplicar\ \ ·\ \ $\delta$ desplegar\ \ ·\ \ $\iota$ recursor sobre constructor\par\vspace{3pt}
\textbf{Proposicional.} $a=b$ es una proposición que hay que \textbf{demostrar}; se usa con \lean{rw}.}
\normal\tarj{e9}{\faLightbulb\ El ejemplo que lo aclara (Ejemplo 10.14)}{
$0+0=0$ sale por \lean{rfl}: se reduce.\\
$0+n=n$ \textbf{no}: $n$ es una variable, no un constructor, y $\iota$ no avanza. Hace falta \lean{induction}.}\grande
\end{minipage}

\vfill
\pruebate{e7}{
\item ¿Qué dice la reducción $\iota$, en una frase?
\item ¿Qué tres reducciones usa el núcleo para comprobar un \lean{rfl}?
\item ¿Por qué $0+n=n$ necesita inducción y $0+0=0$ no?
}{\textbf{1.} Cómo se reduce el recursor aplicado a un constructor.\quad
\textbf{2.} $\beta$, $\delta$ y $\iota$.\quad
\textbf{3.} La suma recurre sobre el segundo argumento: con $0$ se reduce; con la variable $n$, no.}
\navegacion{est6}{Estación 6}{est8}{Estación 8}

\newpage
%=====================================================================
% ESTACIÓN 8: EL CIC DE LEAN
%=====================================================================
\grande
\estacion{8}{\faMicrochip}{El CIC de Lean y las tácticas}{dos precisiones que te van a pedir:
Lean no es «el CIC» tal cual, y una táctica no es una prueba.}

\begin{minipage}[t]{0.47\textwidth}\vspace{0pt}
\tarj{e8}{\faNotEqual\ Tres diferencias con el CIC de Coq (Nota 10.7)}{
\textbf{1. Irrelevancia de pruebas por definición.} Dos pruebas de la misma proposición son
iguales por una \textbf{regla}, no por un axioma.\par\vspace{3pt}
\textbf{2. Cocientes primitivos.} El núcleo trae \lean{Quot} con su regla de cómputo; solo
\lean{Quot.sound} es axioma.\par\vspace{3pt}
\textbf{3. Universos no acumulativos.} Un tipo de \lean{Type 0} no habita automáticamente \lean{Type 1};
se sube con \lean{ULift}. A cambio hay unicidad de tipos.\par\vspace{2pt}
{\normalsize Fuentes: Carneiro (2019); de Moura y Ullrich (2021).}}
\normal\tarj{rojo}{\faPlusCircle\ Los axiomas van encima, no dentro}{
\lean{propext}, \lean{Classical.choice} y \lean{Quot.sound} \textbf{no} son reglas del CIC: se añaden, y
\lean{\#print axioms} los lista.\\[2pt]
\lean{completitud\_ND}: los tres.\\
\lean{quicksort\_correcto}: \lean{propext} y \lean{Quot.sound}.}\grande
\end{minipage}\hfill
\begin{minipage}[t]{0.51\textwidth}\vspace{0pt}
\begin{tcolorbox}[colback=e8!5, colframe=e8, boxrule=1.2pt, arc=5pt, left=5pt, right=5pt,
  title={\faTools\ Cada táctica construye un término del CIC (Tabla 10.2)},
  colbacktitle=e8, fonttitle=\bfseries\large]
\renewcommand{\arraystretch}{1.22}\normalsize
\begin{tabular}{@{}lll@{}}
\textbf{Táctica} & \textbf{Construye} & \textbf{Regla lógica}\\\hline
\lean{intro} & una $\lambda$:\ \ \lean{fun h => …} & introducción de $\to$ y $\forall$\\
\lean{apply}, \lean{exact} & una aplicación:\ \ \lean{h a} & modus ponens\\
\lean{constructor} & \lean{And.intro} & introducción de $\land$\\
\lean{left}, \lean{right} & \lean{Or.inl}, \lean{Or.inr} & introducción de $\lor$\\
\lean{cases} & el eliminador \lean{.casesOn} & eliminación por casos\\
\lean{induction} & el recursor \lean{.rec} & inducción\\
\lean{rfl} & reducir y comparar & igualdad por definición
\end{tabular}
\end{tcolorbox}
\normal\tarj{e8}{\faCodeBranch\ \lean{cases} o \lean{induction}}{
Las dos abren \textbf{un caso por constructor}. Solo \lean{induction} entrega una
\textbf{hipótesis inductiva} por cada argumento recursivo.}
\clave{e8}{Una táctica es un \textbf{metaprograma} que escribe un término. El núcleo comprueba el
término final, no la táctica.}
\end{minipage}

\vfill
\pruebate{e8}{
\item Di las tres diferencias entre el CIC de Lean y el de Coq.
\item ¿Qué término construye \lean{intro}? ¿Y \lean{induction}?
\item ¿Los tres axiomas de la auditoría son parte del CIC?
}{\textbf{1.} Irrelevancia de pruebas por definición, cocientes primitivos, universos no acumulativos.\quad
\textbf{2.} Una $\lambda$; una aplicación del recursor \lean{.rec}.\quad
\textbf{3.} No: se añaden encima, y por eso se auditan.}
\navegacion{est7}{Estación 7}{est9}{Estación 9}

\newpage
%=====================================================================
% ESTACIÓN 9: EL CIC EN MI TESIS
%=====================================================================
\grande
\estacion{9}{\faGraduationCap}{El CIC en mi tesis}{aquí se cierra el círculo. Señala tu código y di
qué pieza del CIC es cada línea.}

\begin{minipage}[t]{0.53\textwidth}\vspace{0pt}
\codigo{e9}{\faFont\ \texttt{Formula}: un dato\ \ ·\ \ \texttt{ND}: una proposición, un constructor por regla}{
inductive Formula : Type\\
| atom : String → Formula\\
| neg\ \ : Formula → Formula\\
| impl : Formula → Formula → Formula\\[6pt]
inductive ND : Set Formula → Formula → Prop\\
| hyp\ \ \ \ \ \ \ : φ ∈ Γ → ND Γ φ\\
| impI\ \ \ \ \ \ : ND (insert φ Γ) ψ → ND Γ (φ $\longrightarrow$ ψ)\\
| impE\ \ \ \ \ \ : ND Γ (φ $\longrightarrow$ ψ) → ND Γ φ → ND Γ ψ\\
| negI\ \ \ \ \ \ : ND (insert φ Γ) Falsedad → ND Γ (\textasciitilde φ)\\
| negE\ \ \ \ \ \ : ND Γ (\textasciitilde φ) → ND Γ φ → ND Γ Falsedad\\
| botE\ \ \ \ \ \ : ND Γ Falsedad → ND Γ φ\\
| classical : ND (insert (\textasciitilde φ) Γ) Falsedad → ND Γ φ\\[2pt]
{\rmfamily\footnotesize\color{black!60}Capítulo 11. Se omiten los argumentos implícitos \{Γ φ ψ\} de cada constructor.}}
\normal\tarj{e7}{\faSortAmountDown\ Quicksort (Capítulo 12)}{
\lean{List} es inductivo (\lean{nil}, \lean{cons}), pero quicksort \textbf{no} es recursión estructural:
llama sobre listas filtradas, no sobre la cola. Por eso lleva \lean{termination\_by l => l.length} y
una prueba de que la longitud decrece.}\grande
\end{minipage}\hfill
\begin{minipage}[t]{0.455\textwidth}\vspace{0pt}
\normal\tarj{e9}{\faLayerGroup\ Los cuatro tiempos, en \lean{ND}}{
\renewcommand{\arraystretch}{1.2}
\begin{tabular}{@{}l>{\raggedright\arraybackslash}p{8.0cm}@{}}
\textbf{Formación} & \lean{ND : Set Formula → Formula → Prop}\\
\textbf{Introducción} & 7 constructores $=$ 7 reglas de inferencia\\
\textbf{Eliminación} & \lean{ND.rec}: inducción sobre la derivación\\
\textbf{Cómputo} & no hace falta: en \lean{Prop} dos pruebas ya son iguales
\end{tabular}}\grande
\tarj{e6}{\faStream\ «Por inducción sobre la derivación»}{
\lean{induction d} aplica \lean{ND.rec}: \textbf{7 casos}, uno por regla, con una hipótesis inductiva
por cada premisa que es una derivación. Así pruebo el \textbf{debilitamiento} y la \textbf{corrección}.\par\vspace{2pt}
\lean{induction φ} aplica \lean{Formula.rec}: \textbf{3 casos}.}
\clave{e9}{Un teorema es un término. \lean{completitud\_ND} tiene tipo
\lean{EsTautologia φ → ND ∅ φ}, y el núcleo comprueba ese tipo.}
\end{minipage}

\vfill
\pruebate{e9}{
\item ¿Cuántos casos tiene una inducción sobre \lean{ND}? ¿Y sobre \lean{Formula}?
\item ¿Qué pieza del CIC es la regla \lean{impI}?
\item ¿Por qué quicksort necesita \lean{termination\_by}?
}{\textbf{1.} Siete y tres: uno por constructor.\quad
\textbf{2.} Un constructor de \lean{ND}: una regla de introducción del tipo inductivo.\quad
\textbf{3.} Porque no llama sobre una componente inmediata del argumento; hay que dar una medida que decrece.}
\navegacion{est8}{Estación 8}{preguntas}{Preguntas}

\newpage
%=====================================================================
% PREGUNTAS 1
%=====================================================================
\pagina{preguntas}{rojo}{\faQuestionCircle}{Preguntas imprescindibles (1 de 3): qué es el CIC}{tapa
la respuesta, contesta en voz alta y compara. Las de \faStar\ son las que no pueden fallar.}

\begin{tcbraster}[raster columns=2, raster column skip=8pt, raster row skip=5pt,
  raster equal height=rows]
\preg*{e2}{1}{¿Qué es el CIC, en una frase?}{Es el Cálculo de Construcciones extendido con tipos
inductivos: tipos que se declaran por sus constructores y que traen su eliminador y sus reglas de
cómputo $\iota$. Con variantes, es la teoría de tipos que implementa el núcleo de Lean.}{estaciones \chip{2} y \chip{8} · \ir[e2]{guion}{\faMicrophone\ Guion}}
\preg*{e1}{2}{¿Qué le falta al CoC?}{Tipos de datos con inducción. En el CoC puro los naturales se
codifican como funciones, los numerales de Church. Eso da iteración, pero el principio de
inducción dependiente no se obtiene de la codificación: lo cito de Geuvers (2001).}{estación \chip{1} · §10.2.3}
\preg*{e2}{3}{¿Qué genera una declaración inductiva?}{Cuatro piezas. La formación: el tipo. La
introducción: los constructores, que son las únicas maneras de construir. La eliminación: el
recursor, la única manera de usar. Y el cómputo: la regla $\iota$.}{estación \chip{2} · Nota 10.6}
\preg*{e6}{4}{¿Por qué el recursor es recursión e inducción?}{Porque es un solo término con un
tipo de retorno variable, el motivo. Si el motivo no depende del valor, define funciones: es la
recursión primitiva. Si depende del valor y vive en \lean{Prop}, su tipo es el principio de
inducción.}{estación \chip{6} · §10.3.1}
\preg*{e7}{5}{¿Qué es la reducción $\iota$?}{La regla de cómputo del eliminador. Aplicado a un
constructor, elige el manejador de ese constructor, le pasa sus argumentos y repite en los
argumentos recursivos. Para $\N$: en cero devuelve la base; en $S\,n$ aplica el paso a $n$ y al
resultado recursivo.}{estación \chip{7} · §10.3.5}
\preg*{e4}{6}{¿Qué es la positividad estricta?}{El tipo que se define no puede
aparecer a la izquierda de una flecha en los argumentos de sus constructores, a ninguna
profundidad. Se exige porque, con un constructor como $(X\to\bot)\to X$, se construiría un
habitante de $\bot$: el sistema sería inconsistente.}{estación \chip{4} · §10.3.2}
\end{tcbraster}
\navegacion{est9}{Estación 9}{preguntas2}{Preguntas 2}

\newpage
%=====================================================================
% PREGUNTAS 2
%=====================================================================
\pagina{preguntas2}{rojo}{\faQuestionCircle}{Preguntas imprescindibles (2 de 3): sorts, Lean y tácticas}{aquí
te piden precisión. Distingue siempre regla, axioma y táctica.}

\begin{tcbraster}[raster columns=2, raster column skip=8pt, raster row skip=5pt,
  raster equal height=rows]
\preg*{e3}{7}{¿Qué diferencia hay entre \lean{Prop} y \lean{Type}?}{Son sorts, tipos de tipos. \lean{Prop}
contiene proposiciones: sus pruebas son irrelevantes y se pueden borrar. \lean{Type} contiene
datos, que sí se distinguen y se computan. En mi tesis, \lean{Formula} vive en \lean{Type} porque
es sintaxis, y \lean{ND} en \lean{Prop} porque es una afirmación.}{estación \chip{3} · Nota 10.2}
\preg{e3}{8}{¿Qué significa que \lean{Prop} sea impredicativa?}{Que una proposición puede cuantificar
sobre todas las proposiciones, incluida ella misma, y sigue siendo una proposición:
$\forall p{:}\texttt{Prop},\ p$ vive en \lean{Prop}. La jerarquía \lean{Type u} es predicativa:
cuantificar sobre un universo sube de universo.}{estación \chip{3}}
\preg{e6}{9}{¿Por qué un inductivo de \lean{Prop} solo se elimina hacia \lean{Prop}?}{Por dos razones
que da la tesis: la eliminación fuerte sobre una sort impredicativa haría inconsistente el sistema,
y las pruebas dejarían de ser borrables. Las excepciones son el caso vacío y el unitario. Por eso
el motivo de \lean{ND.rec} solo llega a \lean{Prop}.}{estación \chip{6} · Nota 10.5}
\preg*{e8}{10}{¿Lean implementa exactamente el CIC?}{No: una variante. Tres diferencias. La
irrelevancia de pruebas vale por definición. Los cocientes son primitivos y solo
\lean{Quot.sound} es axioma. Y los universos no son acumulativos. La descripción precisa es de
Carneiro (2019).}{estación \chip{8} · Nota 10.7}
\preg*{e8}{11}{¿Los axiomas de su auditoría son parte del CIC?}{No. Las reglas del CIC no los
incluyen: \lean{propext}, \lean{Classical.choice} y \lean{Quot.sound} se añaden encima, y
\lean{\#print axioms} los lista. La completitud usa los tres; la corrección de quicksort, solo
\lean{propext} y \lean{Quot.sound}.}{estación \chip{8}}
\preg*{e8}{12}{¿Qué relación hay entre las tácticas y el CIC?}{Una táctica es un metaprograma que
construye un término del CIC. \lean{intro} construye una $\lambda$; \lean{apply}, una aplicación;
\lean{constructor}, un constructor; \lean{cases} e \lean{induction}, el eliminador. El núcleo
comprueba el término final, no las tácticas.}{estación \chip{8} · Tabla 10.2}
\end{tcbraster}
\navegacion{preguntas}{Preguntas 1}{preguntas3}{Preguntas 3}

\newpage
%=====================================================================
% PREGUNTAS 3
%=====================================================================
\pagina{preguntas3}{rojo}{\faQuestionCircle}{Preguntas imprescindibles (3 de 3): tu tesis}{termina
cada respuesta volviendo a lo que hiciste. Y di sin dudar qué es tuyo y qué citas.}

\begin{tcbraster}[raster columns=2, raster column skip=8pt, raster row skip=5pt,
  raster equal height=rows]
\preg*{e9}{13}{¿Dónde usa su tesis el CIC?}{En las definiciones y en cada inducción. \lean{Formula}
es un tipo inductivo con tres constructores, y \lean{ND} es una familia inductiva en \lean{Prop} con
siete, uno por regla. «Inducción sobre la derivación» es aplicar \lean{ND.rec}: siete casos. Así
pruebo la corrección y el debilitamiento.}{estación \chip{9}}
\preg*{gris}{14}{¿Qué demostró usted sobre el CIC?}{Nada: el Capítulo 10 es exposición. Las reglas son
las de Paulin-Mohring (1993), la metateoría es de Werner (1994) y la teoría de Lean la describe
Carneiro (2019). Mi aporte está en los Capítulos 11 y 12, donde uso el CIC para formalizar.}{\ir[e2]{guion}{\faMicrophone\ Guion}}
\preg{e8}{15}{¿\lean{cases} o \lean{induction}?}{Las dos analizan por casos: un caso por constructor.
Solo \lean{induction} entrega una hipótesis inductiva por cada argumento recursivo, porque usa
\lean{.rec}; \lean{cases} usa \lean{.casesOn}. Si necesito $P(k)$ para probar $P(S\,k)$, necesito
\lean{induction}.}{estación \chip{8} · §10.4.1}
\preg{e7}{16}{¿Igualdad por definición o igualdad proposicional?}{Por definición: el núcleo reduce
los dos lados con $\beta$, $\delta$ e $\iota$ y compara; se cierra con \lean{rfl}. Proposicional:
$a=b$ es una proposición que hay que demostrar, y se usa con \lean{rw}. $0+0=0$ sale por
\lean{rfl}; $0+n=n$ necesita inducción.}{estación \chip{7}}
\preg{e9}{17}{Quicksort no es recursión estructural. ¿Cómo se acepta?}{Las llamadas son sobre listas
filtradas, no sobre la cola. Doy una medida, la longitud, y pruebo que decrece en cada llamada:
\lean{termination\_by} y \lean{decreasing\_by}. Como $(\N,<)$ es bien fundado, la función es
total.}{estación \chip{9} · Capítulo 12}
\preg{gris}{18}{¿Cómo sabe que el sistema es consistente?}{No lo demuestro. La tesis cita a Werner
para la metateoría del CIC y a Carneiro para la teoría de Lean. {\small(Fuera del texto de la tesis:
Carneiro da un modelo en teoría de conjuntos, así que la consistencia es relativa.)} Lo que
garantizo es que mis pruebas pasan el núcleo y que la auditoría lista los axiomas.}{estación \chip{8}}
\end{tcbraster}
\navegacion{preguntas2}{Preguntas 2}{trampas}{Trampas}

\newpage
%=====================================================================
% TRAMPAS Y CÓMO DECIRLO
%=====================================================================
\pagina{trampas}{e4}{\faBomb}{Diez trampas y cómo leer los símbolos}{frases que suenan
bien y son falsas. Si te oyes decir una, corrígete en el momento.}

\newcommand{\trampa}[3]{\begin{tcolorbox}[colback=white, colframe=black!35, boxrule=0.8pt, arc=4pt,
  left=6pt, right=6pt, top=2pt, bottom=2pt, before upper=\raggedright]
  \medio\chip{#3}\ {\color{rojo}\faTimes}\ {\color{rojo!80!black}«#1»}\\
  {\color{verde!70!black}\faCheck}\ #2\end{tcolorbox}}

\begin{minipage}[t]{0.62\textwidth}\vspace{0pt}
\begin{tcbraster}[raster columns=2, raster column skip=6pt, raster row skip=4pt,
  raster equal height=rows]
\trampa{El CIC es el CoC.}{Es el CoC \textbf{más} tipos inductivos, con eliminadores y $\iota$.}{2}
\trampa{Con Church no hay recursión.}{Sí hay iteración. Lo que falta es la inducción dependiente.}{1}
\trampa{Lean es el CIC.}{Es una variante: tres diferencias (Nota 10.7).}{8}
\trampa{Ind, Constr y Elim son comandos de Lean.}{Son pseudo-términos de Paulin-Mohring.}{5}
\trampa{Positiva y estrictamente positiva es lo mismo.}{En $(X\to A)\to B$ es covariante y no estricta.}{4}
\trampa{\lean{Prop} es \lean{Bool}.}{\lean{Prop} es un sort de proposiciones; \lean{Bool} es un dato con dos valores.}{3}
\trampa{Los axiomas son parte del CIC.}{Se añaden encima; por eso se auditan.}{8}
\trampa{\lean{rfl} prueba cualquier igualdad.}{Solo las que valen por definición.}{7}
\trampa{La táctica es la prueba.}{La prueba es el término; la táctica lo escribe.}{8}
\trampa{Demostré propiedades del CIC.}{No: lo expongo y lo uso. Cito a Paulin-Mohring y a Werner.}{9}
\end{tcbraster}
\end{minipage}\hfill
\begin{minipage}[t]{0.36\textwidth}\vspace{0pt}
\begin{tcolorbox}[colback=e4!6, colframe=e4, boxrule=1.1pt, arc=5pt, left=6pt, right=6pt,
  top=4pt, bottom=4pt, title={\faVolumeUp\ Cómo se lee en voz alta}, colbacktitle=e4,
  fonttitle=\bfseries\large]
\renewcommand{\arraystretch}{1.4}\medio
\begin{tabular}{@{}l>{\raggedright\arraybackslash}p{5.6cm}@{}}
$(x{:}A)\,B$ & «para todo $x$ de $A$, $B$»: el producto dependiente\\
$[x{:}A]\,t$ & «lambda $x$ de $A$, $t$»\\
$\Ind(X{:}A)\{\dots\}$ & «el inductivo $X$ de aridad $A$ con constructores\ldots»\\
$\Constr(i,I)$ & «el constructor $i$ de $I$»\\
$\Elim(c,Q)\{f_i\}$ & «eliminar $c$ hacia $Q$ con los manejadores $f_i$»\\
$C_i\{I,Q\}$ & «el tipo del manejador del caso $i$»\\
$\ired$ & «se reduce por iota»\\
\lean{motive} & «el motivo»: lo que se prueba o se calcula\\
\end{tabular}
\end{tcolorbox}
\end{minipage}

\vfill
\begin{center}
{\large\bfseries\color{gris!80!black}\faHistory\ Quién hizo qué: di el nombre y el año sin dudar}\\[5pt]
\begin{tikzpicture}[font=\small, >={Stealth[length=2.6mm]},
  hito/.style={draw=#1, fill=#1!10, line width=1pt, rounded corners=4pt, align=center,
    text width=3.75cm, minimum height=1.55cm, inner sep=3pt}]
  \node[hito=e1] (a) {\textbf{1988}\\ Coquand y Huet\\ el CoC};
  \node[hito=e2, right=3.5mm of a] (b) {\textbf{1990}\\ Coquand y Paulin\\ tipos inductivos};
  \node[hito=e5, right=3.5mm of b] (c) {\textbf{1993}\\ Paulin-Mohring\\ las reglas que sigo};
  \node[hito=e6, right=3.5mm of c] (d) {\textbf{1994}\\ Werner\\ metateoría del CIC};
  \node[hito=e8, right=3.5mm of d] (e) {\textbf{2019}\\ Carneiro\\ la teoría de Lean};
  \node[hito=e8, right=3.5mm of e] (f) {\textbf{2021}\\ de Moura y Ullrich\\ Lean 4};
  \foreach \u/\v in {a/b,b/c,c/d,d/e,e/f} \draw[->, line width=1.1pt, black!55] (\u) -- (\v);
\end{tikzpicture}
\end{center}
\navegacion{preguntas3}{Preguntas 3}{reglas}{Hoja de reglas}

\newpage
%=====================================================================
% HOJA DE REGLAS
%=====================================================================
\pagina{reglas}{gris}{\faScroll}{Hoja de reglas}{todas las reglas del CIC en una página. No hay que
memorizarlas: hay que saber leer cada una. Son de Paulin-Mohring (1993).}

\begin{tcbraster}[raster columns=2, raster column skip=8pt, raster row skip=5pt,
  raster equal height=rows]
\begin{tcolorbox}[colback=e5!5, colframe=e5, boxrule=1.1pt, arc=5pt,
  title={Formación (Ind)\hfill\chip{5}}, colbacktitle=e5, fonttitle=\bfseries]
\[\frac{\Gamma,X:A\vdash C_i:s}{\Gamma\vdash\Ind(X:A)\{C_1\mid\dots\mid C_n\}:A}\]
{\small\centering $\mathrm{arity}(A,s)$,\ \ $\mathrm{constructor}(C_i,X)$,\ \ para todo $i$\par}
\end{tcolorbox}
\begin{tcolorbox}[colback=e5!5, colframe=e5, boxrule=1.1pt, arc=5pt,
  title={Introducción (Constr)\hfill\chip{5}}, colbacktitle=e5, fonttitle=\bfseries]
\[\frac{\Gamma\vdash I:A}{\Gamma\vdash\Constr(i,I):C_i[X:=I]}\]
{\small\centering $I=\Ind(X:A)\{C_1\mid\dots\mid C_n\}$,\ \ $1\le i\le n$\par}
\end{tcolorbox}
\begin{tcolorbox}[colback=e6!5, colframe=e6, boxrule=1.1pt, arc=5pt,
  title={Eliminación no dependiente (Nodep)\hfill\chip{6}}, colbacktitle=e6, fonttitle=\bfseries]
\small\[\frac{\Gamma\vdash c:(I\,\vec a)\quad\Gamma\vdash Q:(\vec x{:}\vec A)\,s'\quad\Gamma\vdash f_i:C_i\{I,Q\}}
{\Gamma\vdash\Elim(c,Q)\{f_1\mid\dots\mid f_n\}:(Q\,\vec a)}\]
\end{tcolorbox}
\begin{tcolorbox}[colback=e6!5, colframe=e6, boxrule=1.1pt, arc=5pt,
  title={Eliminación dependiente (Dep)\hfill\chip{6}}, colbacktitle=e6, fonttitle=\bfseries]
\footnotesize\[\frac{\Gamma\vdash c:(I\,\vec a)\quad\Gamma\vdash Q:(\vec x{:}\vec A)(I\,\vec x)\to s'\quad
\Gamma\vdash f_i:C_i\{I,Q,\Constr(i,I)\}}
{\Gamma\vdash\Elim(c,Q)\{f_1\mid\dots\mid f_n\}:(Q\,\vec a\,c)}\]
\end{tcolorbox}
\begin{tcolorbox}[colback=e7!5, colframe=e7, boxrule=1.1pt, arc=5pt,
  title={Reducción $\iota$\hfill\chip{7}}, colbacktitle=e7, fonttitle=\bfseries]
\small\[\Elim((\Constr(i,I)\ \vec m),Q)\{\vec f\}\ \ired\
\bigl(C_i[I,\ \mathrm{Fun\_Elim}(I,Q,\vec f),\ f_i]\ \vec m\bigr)\]
\[\mathrm{Fun\_Elim}(I,Q,\vec f)\ \equiv\ [\vec x{:}\vec A]\,[c{:}(I\,\vec x)]\ \Elim(c,Q)\{\vec f\}\]
\end{tcolorbox}
\begin{tcolorbox}[colback=e2!5, colframe=e2, boxrule=1.1pt, arc=5pt,
  title={Todo, para $\N$\hfill\chip{2}}, colbacktitle=e2, fonttitle=\bfseries]
\medio\raggedright $\texttt{nat}\equiv\Ind(X:\Set)\{X\mid X\to X\}$\hfill $0:\N$\quad $S:\N\to\N$\\[3pt]
$\Elim(0,Q)\{f_0\mid f_S\}\ired f_0$\\[2pt]
$\Elim(S\,n,Q)\{f_0\mid f_S\}\ired f_S\ n\ (\Elim(n,Q)\{f_0\mid f_S\})$\\[3pt]
{\small Manejadores: $Q$ y $\N\to Q\to Q$;\ \ dependientes: $Q\,0$ y $(n{:}\N)\to Q\,n\to Q\,(S\,n)$}
\end{tcolorbox}
\end{tcbraster}

\begin{tcbraster}[raster columns=2, raster column skip=8pt, raster equal height]
\begin{tcolorbox}[colback=e3!6, colframe=e3, boxrule=1.1pt, arc=5pt,
  title={Sorts y notación\hfill\chip{3}}, colbacktitle=e3, fonttitle=\bfseries]
\medio\raggedright $t ::= s\mid x\mid (x{:}t)\,t\mid [x{:}t]\,t\mid (t\,t)$\hfill $Ar ::= s\mid(x{:}M)\,Ar$\\[3pt]
CoC:\ \ $\mathcal S=\{\Set,\Type_S\}$\quad $\mathcal A=\{\Set:\Type_S\}$\quad $\mathcal R=\mathcal S\times\mathcal S$\\[3pt]
$\Pi x{:}A.\,B\leftrightarrow(x{:}A)\,B$\quad $\lambda x{:}A.\,t\leftrightarrow[x{:}A]\,t$\quad $*\leftrightarrow\Set$\quad $\square\leftrightarrow\Type_S$
\end{tcolorbox}
\begin{tcolorbox}[colback=e6!6, colframe=e6, boxrule=1.1pt, arc=5pt,
  title={El principio de inducción es un tipo\hfill\chip{6}}, colbacktitle=e6, fonttitle=\bfseries]
\medio\centering $\forall P:\N\to\Prop,\ \ P\,0\ \to\ (\forall u,\ P\,u\to P\,(S\,u))\ \to\ \forall n,\ P\,n$\\[3pt]
{\small Restricción: de \lean{Prop} solo hacia \lean{Prop}, salvo vacío y unitario (Nota 10.5)}
\end{tcolorbox}
\end{tcbraster}
\navegacion{trampas}{Trampas}{indice}{Índice}

\newpage
%=====================================================================
% ÍNDICE Y LISTA FINAL
%=====================================================================
\pagina{indice}{gris}{\faFont}{Índice A–Z y lista final}{cada término lleva a su estación.
Si dudas con una palabra, haz clic en su número.}

\newcommand{\idx}[2]{#1\dotfill #2\par}
\begin{minipage}[t]{0.61\textwidth}\vspace{0pt}
\setlength{\columnsep}{14pt}
\begin{multicols}{3}
\raggedright\normalsize\setlength{\parskip}{2.5pt}
\idx{Aridad}{\chip{4}}
\idx{Axiomas de Lean}{\chip{8}}
\idx{\lean{cases}}{\chip{8}}
\idx{Church, numerales}{\chip{1}}
\idx{CoC}{\chip{1}}
\idx{Cocientes, \lean{Quot}}{\chip{8}}
\idx{Codificación impredicativa}{\chip{1}}
\idx{Cómputo}{\chip{7}}
\idx{Constr}{\chip{5}}
\idx{Constructor}{\chip{2}\,\chip{4}}
\idx{Eliminación dependiente}{\chip{6}}
\idx{Eliminación singleton}{\chip{6}}
\idx{Eliminador}{\chip{6}}
\idx{Familia inductiva}{\chip{4}}
\idx{\lean{Formula}}{\chip{9}}
\idx{Formación}{\chip{2}\,\chip{5}}
\idx{GTS, PTS}{\chip{3}}
\idx{Igualdad por definición}{\chip{7}}
\idx{Impredicatividad}{\chip{3}}
\idx{Ind}{\chip{5}}
\idx{Índice y parámetro}{\chip{4}}
\idx{\lean{induction}}{\chip{8}}
\idx{Inducción}{\chip{6}}
\idx{Introducción}{\chip{2}\,\chip{5}}
\idx{$\iota$-reducción}{\chip{7}}
\idx{Irrelevancia de pruebas}{\chip{3}\,\chip{8}}
\idx{Iteración}{\chip{1}}
\idx{Leibniz, igualdad}{\chip{1}}
\idx{Manejador}{\chip{6}}
\idx{Metaprograma}{\chip{8}}
\idx{Motivo}{\chip{6}}
\idx{\lean{ND}}{\chip{9}}
\idx{Positividad estricta}{\chip{4}}
\idx{Predecesor}{\chip{1}}
\idx{\lean{Prop}}{\chip{3}}
\idx{Pseudo-término}{\chip{5}}
\idx{Recursión estructural}{\chip{9}}
\idx{Recursor, \lean{.rec}}{\chip{6}}
\idx{\lean{rfl}}{\chip{7}}
\idx{Sort}{\chip{3}}
\idx{Táctica}{\chip{8}}
\idx{\lean{termination\_by}}{\chip{9}}
\idx{Torniquete $\vdash$}{\chip{5}}
\idx{\lean{Type u}}{\chip{3}}
\idx{Universos}{\chip{8}}
\end{multicols}
\end{minipage}\hfill
\begin{minipage}[t]{0.37\textwidth}\vspace{0pt}
\begin{tcolorbox}[colback=verde!6, colframe=verde, boxrule=1.2pt, arc=6pt,
  title={\faClipboardCheck\ Estoy listo si puedo\ldots}, colbacktitle=verde, fonttitle=\bfseries\large,
  before upper=\raggedright]
\medio\newcommand{\si}{{\Large$\square$}\ }
\si decir el guion de 1 minuto sin leer\\[3pt]
\si explicar qué le falta al CoC\\[3pt]
\si nombrar los cuatro tiempos\\[3pt]
\si escribir \lean{Nat} y el tipo de \lean{Nat.rec}\\[3pt]
\si escribir las dos reglas $\iota$ de $\N$\\[3pt]
\si explicar la positividad estricta\\[3pt]
\si distinguir \lean{Prop} de \lean{Type}\\[3pt]
\si decir las tres diferencias de Lean\\[3pt]
\si señalar el CIC en \lean{Formula} y \lean{ND}\\[3pt]
\si contestar las 12 preguntas con \faStar
\end{tcolorbox}
\end{minipage}

\vfill
\begin{tcolorbox}[colback=tldr, colframe=e9, boxrule=1.2pt, arc=6pt, left=10pt, right=10pt,
  top=5pt, bottom=5pt, title={\faStar\ Si solo recuerdas cinco frases},
  colbacktitle=e9, fonttitle=\bfseries\large]
\large
\chip{1}\ \ \textbf{El problema:} codificar datos con funciones da iteración, no inducción.\\[2pt]
\chip{2}\ \ \textbf{La solución:} declarar el tipo por sus constructores.\\[2pt]
\chip{6}\ \ \textbf{El recursor:} una sola pieza; con datos es recursión, con pruebas es inducción.\\[2pt]
\chip{7}\ \ \textbf{$\iota$:} el recursor frente a un constructor se reduce al manejador de ese constructor.\\[2pt]
\chip{9}\ \ \textbf{Mi tesis:} \lean{Formula} y \lean{ND} son inductivos; mis inducciones son sus recursores.
\hfill\ir[rojo]{preguntas}{\faQuestionCircle\ Repasar las preguntas}
\end{tcolorbox}
\navegacion{reglas}{Hoja de reglas}{portada}{Portada}

\end{document}
